\begin{frame}{DP da moeda}
    \metroset{block=fill}

    \begin{block}{Definição}
        Seja $f$ uma função recursiva. Chamaremos um elemento $i$ do domínio da função de {\bf estado} e chamaremos o código que faz as chamadas recursivas de {\bf transição}.
    \end{block}

    O código da função dos slides anteriores é:

    \begin{block}{Implementação da memoização}
        \footnotesize
        \inputminted{cpp}{codigos/moeda.cpp}    
    \end{block}
    A complexidade de uma DP é a soma da complexidade da transição de cada estado. Como a transição do código acima é sempre constante, então a complexidade é $O(estados) = O(n)$.
\end{frame}
